% !PRINTED 339
\weberpage{339}
%<<CONT

1. \gesperrt{Ist $d\equiv 2, 3\pmod 4$, so sind alle und nur die
Zahlen des K\"orpers $\Omega$ ganz, die in der Form}
\[
x+y\sqrt{d}
\]
\gesperrt{mit ganzem rationalen $x, y$ enthalten sind.}

Ist aber $d\equiv 1\pmod 4$, so ist (3) befriedigt, wenn $x$ und $y$
beide gerade oder beide ungerade sind, also wenn $x$ die Form
$2x_1-y$ hat. Ersetzt man dann wieder $x_1$ durch $x$, so folgt:

2. \gesperrt{Ist $d\equiv 1\pmod 4$, so sind alle und nur die Zahlen
in $\Omega$ ganz, die in der Form}
\[
x+y\,\frac{-1+\sqrt{d}}{2}
\]
\gesperrt{mit ganzen rationalen $x, y$ enthalten sind.}

Es ist also im Falle 1. $(1,\sqrt{d})$, im Falle 2. $\left(1,\dfrac{-1+\sqrt{d}}{2}\right)$
\gesperrt{eine Minimalbasis des K\"orpers $\Omega$. Die Grundzahl oder Diskriminante des K\"orpers $\Omega$} ist also im Falle 1.:
\weq{5}{\varDelta=\begin{vmatrix}1, & \sqrt{d}\\ 1, & -\sqrt{d}\end{vmatrix}^{2}=4d,}
im Falle 2.:
\weq{6}{\varDelta=\begin{vmatrix}1, & \dfrac{-1+\sqrt{d}}{2}\\[6pt] 1, & \dfrac{-1-\sqrt{d}}{2}\end{vmatrix}^{2}=d.}
In beiden F\"allen kann man also die ganze Zahl $\omega$ in die
Form setzen:
\weq{7}{\omega=\frac{x+y\sqrt{\varDelta}}{2},}
mit der Bedingung, da\ss
\weq{8}{4N(\omega)=x^2-\varDelta y^2\equiv 0\pmod 4}
sei.

Die beiden Formen der Minimalbasis, die wir erhalten haben,
n\"amlich $(1,\sqrt{d})$ und $\left(1,\dfrac{-1+\sqrt{d}}{2}\right)$, sind also:
\weq{9}{\begin{aligned}(1,0)&=\left(1,\tfrac{1}{2}\sqrt{\varDelta}\right), & \varDelta&\equiv 0,\\ &=\left(1,\frac{-1+\sqrt{\varDelta}}{2}\right), & \varDelta&\equiv 1.\end{aligned}\quad(\mathrm{mod}\ 4).}
Hierin kann das Vorzeichen von $\sqrt{\varDelta}$ in beliebiger Weise bestimmt werden, soll aber dann in derselben Betrachtung festgehalten werden. Wir wollen ein f\"ur allemal annehmen:
%CONT>>
